\section{Metrics Used in Pose Estimation and Tracking}
Evaluating the performance of pose estimation and tracking models involves a
variety of metrics, each designed to measure the accuracy and effectiveness of
predicted joint positions and overall tracking performance. These metrics
revolve around the accuracy of the predicted points compared to the ground
truth, assessing how well the model identifies and tracks keypoints across
images or video sequences.

\subsection{Pose Estimation Metrics}
Pose estimation metrics primarily focus on the accuracy of detecting individual
body joints. Below are some of the key metrics used in the evaluation of pose
estimation:

\textbf{Percentage of Correct Parts (PCP)} \cite{kiefel2014human}: Initially a
standard in the field, PCP measures the accuracy of limb detection. A limb is
considered correctly detected if the distance between the predicted and actual
joint positions is less than half the limb's length. However, PCP has been
criticized for its bias against shorter limbs, leading to its decline in
popularity.

\textbf{Percentage of Detected Joints (PDJ)} \cite{abobakr2016body}: PDJ was
developed to address the limitations of PCP. It measures the accuracy of joint
detection based on the proximity of the predicted joint to its true position,
relative to the torso's diameter. This metric allows for different levels of
localization precision by adjusting the acceptable error threshold.

\textbf{Percentage of Correct Keypoints (PCK)} \cite{gouidis2019accurate}:
Similar to PDJ, PCK evaluates the proximity of predicted joint locations to
their true counterparts but uses the maximum dimension of the ground truth
bounding box as the reference. This metric provides a measure of localization
accuracy for body joints.

\textbf{PCKh} \cite{andriluka20142d}: An adaptation of PCK, PCKh uses the head
segment length to normalize the matching threshold, accounting for individual
size differences. The threshold is set at 50\% of the head segment length,
promoting a size-independent evaluation of joint detection accuracy.

\textbf{Area Under the Curve (AUC)} \cite{guler2018densepose}: This metric
reviews the model's performance across varying PCK thresholds, illustrating its
capability to accurately localize body joints over a range of precision levels.

\textbf{Object Keypoint Similarity (OKS) \cite{maji2022yolo}}: Analogous to
Intersection over Union (IoU) \cite{rezatofighi2019generalized} in object
detection, OKS assesses how closely the predicted keypoints match the actual
data points, emphasizing the quality of overlap between predicted and true
keypoints.

\subsection{Pose Tracking Metrics}
For multi-person pose tracking, additional metrics are used to evaluate the
accuracy of tracking across frames in a video sequence. These metrics are
derived from Multiple Object Tracking (MOT) metrics \cite{milan2016mot16} and
focus on how well the model maintains the identity and position of detected
joints over time:

\textbf{Multiple Object Tracking Accuracy (MOTA)}
\cite{bernardin2008evaluating}: MOTA is a metric used to evaluate the overall
performance of a tracking system. It takes into account three key factors: false
positives (when the system incorrectly identifies a non-existent object), false
negatives (when the system fails to detect an actual object), and identity
switches (when the system mistakenly assigns a different identity to an already
tracked object). By considering these factors, MOTA provides a comprehensive
assessment of the system's accuracy. It calculates this accuracy by averaging
the tracking performance for each joint across all body joints, making it an
important measure of how effectively the model tracks multiple people over time.

\textbf{Multiple Object Tracking Precision (MOTP)}
\cite{bernardin2008evaluating}: MOTP measures the precision of the predicted
joint positions relative to the ground truth, focusing on how accurately the
model tracks the position of each joint. High MOTP values indicate that the
predicted joint locations are close to their true positions.

\textbf{Precision and Recall}: These metrics evaluate the proportion of
correctly identified joints (Precision) and the proportion of true joints
correctly detected (Recall) (The Average Precision and Average Recall metrics
used in \cite{Chen_2018_CVPR}). High precision and recall indicate that the
model is effective at detecting and tracking body joints without missing or
incorrectly identifying them.

\textbf{Mean Average Precision (mAP)}: Mean Average Precision (mAP) is employed
to assess frame-wise multi-person pose accuracy, as described in
\cite{pishchulin2016deepcut}. The mAP metric evaluates how accurately each
person is detected and localized within a frame, accounting for both detection
and localization errors. Traditionally, mAP evaluation depends on knowing the
location and rough scale of a group of individuals during evaluation. However,
in realistic scenarios, particularly in video contexts, such ground-truth
information is rarely available.

These metrics collectively provide a comprehensive framework for evaluating the
performance of multi-person pose tracking models. By combining measures of
accuracy, precision, and recall with the mAP metric, these evaluations offer
insights into both the localization and tracking capabilities of the models
across complex and dynamic video sequences.

\subsection{Challenges in Multi-Person Pose Tracking}
The evaluation of multi-person pose tracking models reveals that tracking
accuracy significantly decreases with the increased complexity of video
sequences/images. Simple data with clearly separated individuals in upright poses
generally result in high tracking accuracy, while data involving
overlapping individuals, fast movements, and complex body articulations tend to
challenge the tracking algorithms, leading to lower MOTA scores. The
relationship between pose estimation accuracy and tracking performance often
shows that better pose estimation corresponds to better tracking, though there
are exceptions where accurate pose estimation does not translate to effective
tracking due to the challenges posed by complex scenes.